\chapter{Objetivos, Decisões Técnicas e Planeamento}\label{chap3}


O estágio curricular realizado na 20 Recolher --- Gestão de Resíduos, Lda. teve como ponto de partida um levantamento detalhado das necessidades tecnológicas da empresa.

\section{Estado Inicial e Problema Identificado}

No início do período de estágio, verificou-se que uma parte significativa dos processos operacionais dependia de procedimentos manuais e de informação distribuída por diferentes canais. Os pedidos de recolha eram, geralmente, comunicados através de telefone ou correio eletrónico, sem uma plataforma centralizada que permitisse a sua submissão, validação e acompanhamento estruturado.

O registo das pesagens era efetuado manualmente em papel, após a leitura visual dos valores apresentados pela balança industrial. Posteriormente, esse registo era encaminhado para a secretaria, onde a informação era tratada administrativamente. Este processo aumentava a possibilidade de erros de transcrição, perda ou duplicação de informação e dificultava a consulta do histórico das operações.

A presença digital da empresa também não disponibilizava uma área reservada através da qual os clientes pudessem submeter pedidos de recolha, acompanhar o respetivo estado ou consultar documentação associada de forma autónoma. Esta situação resultava numa maior dependência dos contactos telefónicos e das trocas de correio eletrónico.

Neste contexto, o estágio teve como principal finalidade a conceção e o desenvolvimento de um conjunto integrado de soluções tecnológicas destinado a centralizar a informação, apoiar os processos operacionais e melhorar a interação entre a empresa e os seus clientes.

\section{Objetivo Geral}

Conceber e desenvolver um ecossistema integrado de soluções de \textit{software} para apoiar e digitalizar os principais processos da 20 Recolher, desde a submissão de pedidos de recolha por parte dos clientes até ao registo das pesagens realizadas no armazém, promovendo a centralização da informação, a rastreabilidade das operações e a redução da dependência de procedimentos manuais.

\section{Objetivos Específicos}

Para concretizar o objetivo geral, definiram-se os seguintes objetivos específicos:

\begin{itemize}
    \item \textbf{Desenvolvimento do Website Institucional:} Criar um novo website institucional moderno para a empresa, com duas vertentes distintas: uma área pública para apresentação dos serviços e captação de novos clientes, e uma área reservada para clientes registados. Esta área reservada permite que os clientes efetuem pedidos de recolha de forma autónoma, acompanhem o estado das suas recolhas e comuniquem diretamente com os operadores através de um sistema de mensagens diretas. A presença digital estruturada é fundamental para reduzir o volume de contactos telefónicos e e-mails manuais, de forma a centralizar a gestão de clientes numa plataforma digital segura.
    \item \textbf{Desenvolvimento da Aplicação de \textit{Backoffice}:} Construir uma aplicação de gestão interna para os operadores da empresa, que centralize o controlo de clientes, a triagem e agendamento de pedidos de recolha, a visualização de pesagens e a administração de utilizadores e permissões. O objetivo é substituir o uso de folhas de cálculo e comunicações manuais.
    \item \textbf{Desenvolvimento do Terminal de Pesagem (\textit{Kiosk}):} Desenvolver uma aplicação destinada ao apoio das operações de pesagem e concebida para utilização através de um ecrã tátil. A aplicação permite ao operador introduzir manualmente os valores apresentados pela balança industrial e associá-los à recolha e ao respetivo tipo de material, reduzindo a necessidade de registos dispersos e facilitando a centralização da informação.

    \item \textbf{Modelação e Implementação da Base de Dados Centralizada:} Conceber e implementar um modelo de dados relacional capaz de suportar as diferentes aplicações do ecossistema, garantindo a integridade referencial e a centralização da informação. Sempre que aplicável, foi utilizada exclusão lógica, de forma a evitar a eliminação física imediata de registos relevantes e a preservar a sua disponibilidade para consulta.

    \item \textbf{Desenvolvimento da \ac{API} Central:} Criar uma \textit{Application Programming Interface} (API) \textit{RESTful} que funcione como a única camada de acesso à base de dados, garantindo que todas as aplicações comuniquem com os dados de forma centralizada, segura e consistente. Este objetivo surgiu da necessidade de separar as duas aplicações Blazor (backoffice e terminal de pesagem) que inicialmente foram concebidas como uma única aplicação.

    \item \textbf{Testes e Validação Funcional:} Validar o funcionamento das aplicações através de testes manuais realizados em computadores pessoais e num ecrã tátil disponibilizado pela empresa para futura utilização na zona operacional do armazém. Embora o equipamento ainda não se encontrasse instalado no local previsto, foi possível testar a interação por toque e simular o processo de introdução dos valores apresentados pela balança industrial. A integração direta com a balança não foi implementada, pelo que os valores continuaram a ser introduzidos manualmente pelo operador.
\end{itemize}


\section{Decisões Técnicas e Justificação das Escolhas}\label{sec:tecnologias-escolhas}

A escolha das tecnologias baseou-se em critérios como a adequação técnica ao problema, a facilidade de aprendizagem, a estabilidade e o suporte disponível no respetivo ecossistema, bem como a capacidade de integração entre as diferentes componentes do sistema.

\subsection{Desenvolvimento Frontend e Presença Web}
\begin{itemize}
    \item \textbf{\textit{Next.js} 16 (\textit{Website} Institucional):} O \textit{Next.js} foi escolhido por permitir desenvolver, no mesmo projeto, tanto as páginas institucionais como as áreas interativas do portal. Para isso, foram utilizados componentes de servidor e componentes de cliente, consoante as necessidades de cada funcionalidade. Os componentes de servidor são processados no servidor antes de o conteúdo ser enviado para o navegador, sendo adequados para apresentar informação estática, melhorar o desempenho e facilitar a indexação pelos motores de busca. Por sua vez, os componentes de cliente são executados no navegador do utilizador e permitem implementar funcionalidades interativas, como formulários, botões, filtros e atualizações dinâmicas da interface. A arquitetura baseada no \textit{App Router} permitiu ainda organizar o projeto de forma estruturada, através de rotas, páginas e componentes reutilizáveis.
    \item \textbf{\textit{Tailwind \ac{CSS}} v4 (Estilização):} O \textit{Tailwind CSS v4} foi adotado para acelerar o desenvolvimento visual da interface e garantir a sua adaptação a diferentes tamanhos de ecrã, incluindo smartphones, tablets e computadores. Esta tecnologia baseia-se em classes utilitárias, ou seja, pequenas classes previamente definidas que aplicam diretamente propriedades de estilo, como espaçamentos, cores, tamanhos ou alinhamentos. Desta forma, foi possível reduzir a necessidade de criar folhas de estilo extensas e personalizadas, diminuindo a quantidade de código CSS a manter e promovendo uma aparência consistente em todo o website. Como os clientes da empresa são maioritariamente entidades empresariais que podem consultar os seus dados através de vários dispositivos, a responsividade foi considerada um requisito operacional essencial e não apenas uma opção visual.
\end{itemize}

\subsection{Aplicações Empresariais Internas}
\begin{itemize}
    \item \textbf{\textit{Photino.Blazor} com .NET 10 (\textit{Backoffice} e \textit{Kiosk}):} O \textit{Photino.Blazor} foi utilizado no desenvolvimento das aplicações de ambiente de trabalho por permitir apresentar interfaces construídas com Blazor dentro de uma janela nativa do sistema operativo. Desta forma, foi possível desenvolver a interface com C\#, HTML e CSS, mantendo uma base tecnológica semelhante à utilizada na \ac{API} e aproveitando os conhecimentos previamente adquiridos em C\#. A aplicação recorre à \textit{WebView} disponibilizada pelo sistema operativo, isto é, a um componente que permite apresentar conteúdos web dentro de uma aplicação de ambiente de trabalho. Esta abordagem contribuiu para criar aplicações relativamente leves e adequadas aos equipamentos existentes na empresa.
    \item \textbf{\textit{Tauri} (Ferramenta de Toners):} O \textit{Tauri} foi utilizado para transformar a ferramenta de comparação de preços de toners numa aplicação de ambiente de trabalho destinada a uso interno. Esta tecnologia permite executar uma interface desenvolvida com tecnologias web dentro de uma janela nativa, recorrendo também à \textit{WebView} do sistema operativo. Esta escolha adequou-se ao objetivo de disponibilizar uma aplicação simples, leve e com reduzido consumo de recursos, que pudesse ser executada localmente no computador da empresa sem depender de um servidor web dedicado.
\end{itemize}

\subsection{Persistência e Infraestrutura de Dados}
\begin{itemize}
    \item \textbf{\textit{PostgreSQL} com \textit{Supabase}:} O PostgreSQL foi escolhido pela sua fiabilidade e pelo suporte a transações ACID, sigla correspondente a \textit{Atomicidade}, \textit{Consistência}, \textit{Isolamento} e \textit{Durabilidade}. Estas propriedades garantem que as operações são concluídas integralmente, respeitam as regras da base de dados, não interferem incorretamente entre si quando executadas em simultâneo e permanecem guardadas após a confirmação. Por exemplo, quando vários operadores registam pesagens ao mesmo tempo, o PostgreSQL controla a concorrência entre as transações, evitando registos incompletos ou alterações contraditórias. A plataforma \textit{Supabase} permitiu utilizar PostgreSQL gerido na nuvem, autenticação integrada e armazenamento de ficheiros, reduzindo a necessidade de configurar e manter servidores próprios.
    \item \textbf{Base de Dados Centralizada e \ac{API} Central:} O acesso aos dados foi concentrado numa \ac{API} desenvolvida em .NET 10. Esta camada recebe os pedidos das aplicações, valida os dados e as permissões dos utilizadores, aplica as regras de negócio e comunica com a base de dados. Desta forma, evitou-se a duplicação de lógica e melhorou-se o controlo e a manutenção do sistema.
    \item \textbf{\textit{Entity Framework Core} (\ac{ORM}):} O \textit{Entity Framework Core} foi utilizado como \ac{ORM}, sigla de \textit{Object-Relational Mapping}. Um ORM permite representar tabelas da base de dados através de classes em C\#, relacionando propriedades dos objetos com colunas e registos. Assim, foi possível consultar, adicionar e atualizar dados através de código C\#, sendo o \textit{Entity Framework Core} responsável por gerar os comandos SQL necessários. Foi utilizada uma abordagem \textit{Code-First}, em que o modelo da base de dados é definido no código e atualizado através de migrações, permitindo controlar e reproduzir as alterações do esquema nos diferentes ambientes do projeto.
\end{itemize}

\subsection{Fluxo de Trabalho e Publicação}
\begin{itemize}
    \item \textbf{\textit{Git} e \textit{GitHub}:} O controlo de versões com Git e o alojamento no GitHub foram adotados desde o primeiro dia de desenvolvimento. Esta decisão justifica-se pela necessidade de manter um histórico completo de todas as alterações ao código, possibilitar a reversão segura de mudanças problemáticas e garantir que o projeto não dependia de uma única máquina. Além disso, o Git permitiu que o desenvolvimento de novas funcionalidades ocorresse em ramos (\textit{branches}) isolados, sem perturbar a versão estável em produção.
    \item \textbf{\textit{Vercel} (Publicação Contínua do Website):} A plataforma Vercel foi escolhida para alojar o website institucional durante o desenvolvimento e os testes. A sua integração direta com o GitHub permitiu gerar automaticamente uma nova versão publicada após as atualizações enviadas para o repositório, tornando o website acessível sem necessidade de configurar ou manter servidores próprios. Esta abordagem facilitou a realização de testes em dispositivos reais, como computadores e smartphones, e permitiu verificar a responsividade e o funcionamento da interface ao longo do desenvolvimento.
    \item \textbf{Autenticação com \textit{Supabase Auth}:} O serviço \textit{Supabase Auth} foi utilizado para gerir o registo, a autenticação e as sessões dos utilizadores do portal. Esta opção permitiu recorrer a um serviço especializado, evitando o desenvolvimento integral de um mecanismo próprio de autenticação. O acesso às funcionalidades reservadas foi condicionado pela validação da sessão e das permissões associadas a cada utilizador.
\end{itemize}

\subsection{Validação Funcional e de Integração}
\begin{itemize}
    \item \textbf{Testes em Ambiente Real:} Validação da usabilidade e fluxo de trabalho do terminal quiosque tátil sob condições reais de armazém.
\end{itemize}


\section{Plano de Trabalhos}\label{sec:plano-trabalhos}

O plano de trabalhos foi estruturado em cinco fases, que se sobrepõem temporalmente de forma natural. As fases de análise e modelação ocorreram essencialmente no início do estágio, enquanto o desenvolvimento e os testes decorreram de forma contínua ao longo do período. O estágio decorreu entre 4 de março e 14 de julho de 2026.

\subsection{Fase 1: Levantamento de Requisitos e Análise}

Nesta fase inicial, procedeu-se à análise das necessidades operacionais da empresa através da observação direta dos processos existentes e das dificuldades identificadas no trabalho diário. O objetivo foi compreender o funcionamento da gestão de pedidos, do registo de pesagens e da comunicação com os clientes, identificando as necessidades que poderiam ser apoiadas através das soluções tecnológicas desenvolvidas.

Foram identificados e documentados requisitos de dois tipos:
\begin{itemize}
    \item \textbf{Requisitos Funcionais:} As funcionalidades que o sistema deveria disponibilizar, como a submissão de pedidos de recolha online, o registo manual assistido de pesagens, o sistema de mensagens entre clientes e operadores e a gestão de permissões por cargo;
    \item \textbf{Requisitos Não Funcionais:} As características de qualidade do sistema, como a segurança dos dados, a compatibilidade com dispositivos táteis, o tempo de resposta da \ac{API} e a disponibilidade numa infraestrutura alojada na nuvem.
\end{itemize}

\subsection{Fase 2: Desenho da Arquitetura e Modelação da Base de Dados}

Com base nos requisitos identificados, procedeu-se ao desenho da arquitetura global do sistema. A solução foi organizada em diferentes componentes, incluindo uma \ac{API} central, uma base de dados relacional, aplicações de ambiente de trabalho e um \textit{website}. A \ac{API} foi definida como a principal camada de comunicação com a base de dados, centralizando as operações de acesso e persistência da informação.

Paralelamente, foi desenhado o modelo de dados relacional, identificando as entidades principais, as suas relações de cardinalidade e as restrições de integridade referencial. O modelo foi implementado em \textit{PostgreSQL} através de uma abordagem \textit{Code-First} com \textit{Entity Framework Core}, sendo o esquema gerado e versionado através do mecanismo de migrações.

\subsection{Fase 3: Desenvolvimento e Implementação}

Esta foi a fase mais longa e correspondeu à implementação prática das soluções. O desenvolvimento foi realizado em entreajuda entre os dois estagiários da empresa, com a contínua revisão e alinhamento de decisões técnicas ao longo do processo. As componentes desenvolvidas abrangem:
\begin{itemize}
    \item \textbf{\textit{Website} Institucional e Portal de Cliente:} Desenvolvimento em \textit{Next.js} 16 e \textit{Tailwind \ac{CSS}} v4, com integração com a \ac{API} central;
    \item \textbf{Terminal de Pesagens (\textit{Kiosk}):} Desenvolvimento em \textit{Photino.Blazor} e .NET 10, com interface adaptada à utilização tátil e autenticação simplificada através de PIN;
    \item \textbf{\textit{Backoffice}:} Desenvolvimento da aplicação interna para gestão de clientes, pedidos, recolhas, pesagens e utilizadores;
    \item \textbf{\ac{API} Central:} Implementação das operações de acesso aos dados, autenticação, validação e comunicação entre as diferentes aplicações;
    \item \textbf{Ferramenta de Comparação de Toners:} Desenvolvimento de uma aplicação interna, posteriormente disponibilizada em formato de aplicação de ambiente de trabalho através do \textit{Tauri}.
\end{itemize}


\subsection{Fase 4: Testes e Validação}

Em paralelo com o desenvolvimento, foram realizados testes manuais e funcionais destinados a verificar os principais fluxos das aplicações e a comunicação entre as diferentes componentes do sistema.

A verificação abrangeu duas vertentes:
\begin{itemize}
    \item \textbf{Testes de Interface e Usabilidade:} Verificação da navegação, da introdução de dados e da adaptação da interface do terminal a um ecrã tátil disponibilizado pela empresa;
    \item \textbf{Testes de Integração de Dados:} Confirmação de que os dados submetidos através das diferentes aplicações eram corretamente enviados para a \ac{API} e persistidos na base de dados.
\end{itemize}

Estes testes permitiram identificar e corrigir problemas funcionais e de apresentação ao longo do processo de desenvolvimento. Não foram realizados testes formais com utilizadores finais, constituindo esta uma limitação do trabalho desenvolvido.


\subsection{Fase 5: Documentação e Entrega}

A fase final correspondeu à consolidação da documentação técnica e à preparação da entrega do projeto. Foram produzidos documentos de referência para futuros programadores, nomeadamente um guia de instalação em produção, um guia do \textit{frontend} para clientes e a documentação dos controladores da \ac{API}. Esta fase incluiu igualmente a elaboração do presente relatório de estágio.

\vspace{1em}
O cronograma das fases de trabalho ao longo do período de estágio encontra-se na Tabela~\ref{tab:cronograma}.

\begin{table}[H]
    \centering
    \caption{Cronograma das Fases do Plano de Trabalhos (2026).}
    \label{tab:cronograma}
    \begin{tabularx}{\textwidth}{|X|c|c|c|c|c|}
        \hline
        \rowcolor{gray!20} \textbf{Fase / Atividade} & \textbf{Mar.} & \textbf{Abr.} & \textbf{Mai.} & \textbf{Jun.} & \textbf{Jul.} \\ \hline
        Fase 1: Levantamento de Requisitos e Análise & \cellcolor{gray!50} X & & & & \\ \hline
        Fase 2: Arquitetura e Modelação da Base de Dados & \cellcolor{gray!50} X & \cellcolor{gray!50} X & & & \\ \hline
        Fase 3: Desenvolvimento e Implementação & & \cellcolor{gray!50} X & \cellcolor{gray!50} X & \cellcolor{gray!50} X & \cellcolor{gray!50} X \\ \hline
        Fase 4: Testes e Validação & & \cellcolor{gray!50} X & \cellcolor{gray!50} X & \cellcolor{gray!50} X & \cellcolor{gray!50} X \\ \hline
        Fase 5: Documentação e Entrega & & & & \cellcolor{gray!50} X & \cellcolor{gray!50} X \\ \hline
    \end{tabularx}
\end{table}